\metatitle{Introduction to colored matroids}
\metadescription{Fixed-quota colorings, disagreement cycles, and the symmetric cycle-exchange axiom.}

\section[intro]{Words with fixed quotas}

Let $E$ be a finite set of positions, let $[k]=\{1,\dots,k\}$ be a set of
colors, and fix nonnegative integers $r_1,\dots,r_k$ whose sum is $|E|$.
A \defin{word} is a map $u:E\to[k]$. We write
\[
 B(u,c)=u^{-1}(c)
\]
for the positions carrying color $c$. A \defin{fixed-quota coloring family} is
a nonempty collection $\mathcal F$ of words such that
\[
 |B(u,c)|=r_c
 \qquad\text{for every }u\in\mathcal F\text{ and }c\in[k].
\]
When $E=[n]$, we display a word by listing its colors in position order. Thus
$\mathtt{ABAC}$ denotes the map with values $A,B,A,C$ at positions
$1,2,3,4$.

Ordinary matroid bases are recovered by looking at one color fiber. The
colored object retains all fibers simultaneously, so compatible exchanges
must remember how several colors move around a cycle.

\section[disagreementGraph]{The disagreement graph}

Fix two words $u,v\in\mathcal F$. Their \defin{disagreement graph}
$D(u,v)$ is the directed multigraph whose vertices are the colors. For each
position $e$ at which $u(e)\ne v(e)$, it has an arc
\[
 u(e)\longrightarrow v(e)
\]
labeled by $e$. The quota conditions imply that every vertex has equal
indegree and outdegree.

A \defin{simple color cycle} is a directed cycle in $D(u,v)$ that visits each
of its colors once. Switching its positions means replacing the entries of
one word by those of the other there. Entries outside the cycle are left
unchanged. We denote the two resulting words by $u'$ and $v'$.

Every color occurs equally often in $u'$ and $u$ (and in $v'$ and $v$), so a
switch automatically preserves the quotas. The real condition is that both
resulting words remain in the chosen family $\mathcal F$.

\section[symmetricCycleExchange]{Symmetric cycle exchange}

\begin{definition}[Symmetric cycle exchange]
An \defin{SCE family} is a fixed-quota coloring family $\mathcal F$ satisfying
\defin{symmetric cycle exchange}: for every $u,v\in\mathcal F$ and every
position $e$ with $u(e)\ne v(e)$, the arc of $D(u,v)$ labeled by $e$ lies on a
simple color cycle such that switching its positions in both words gives two
words in $\mathcal F$.
\end{definition}

The position $e$ is prescribed: it is not enough that some switch works. The
condition is also symmetric: both resulting words must belong to $\mathcal F$.

With three colors, a simple color cycle has length two or three. Thus an
exchange may be an ordinary swap between two colors, or it may move three
colors cyclically.

\section[exchangeExamples]{Examples of colored exchange}

\begin{example}[A two-color exchange]
Let $\mathcal F$ be the family of all words on four positions with two $A$'s and two
$B$'s. Take
\[
 u=\mathtt{AABB},\qquad v=\mathtt{BBAA}.
\]
At positions $1$ and $3$, the disagreement arcs are respectively
$A\to B$ and $B\to A$. Switching this two-cycle gives
\[
 u'=\mathtt{BAAB},\qquad v'=\mathtt{ABBA}.
\]
Both resulting words have the required quotas. More generally, with two colors
in the full fixed-quota coloring family, every prescribed disagreement can be
paired with a reverse disagreement. This is the two-color form of symmetric
exchange.
\end{example}

\begin{example}[A three-color cycle]
Consider
\[
 \mathcal F=\{\mathtt{ABC},\mathtt{BCA},\mathtt{CAB}\}.
\]
For $u=\mathtt{ABC}$ and $v=\mathtt{BCA}$, the three positions label the
cycle
\[
 A\longrightarrow B\longrightarrow C\longrightarrow A.
\]
Switching the cycle interchanges $u$ and $v$. The same happens for either of
the other pairs, so $\mathcal F$ satisfies symmetric cycle exchange.
\end{example}

\begin{example}[A switch inside a larger family]
Consider the fixed-quota coloring family
\[
 \mathcal F=\{\mathtt{AABC},\mathtt{ABAC},\mathtt{ACBA},
 \mathtt{BAAC},\mathtt{CABA},\mathtt{BCAA},\mathtt{CBAA}\}.
\]
Take $u=\mathtt{BCAA}$ and $v=\mathtt{AABC}$. The arcs at positions $1$ and
$3$ form the two-cycle $B\to A\to B$. Switching those positions gives
\[
 u'=\mathtt{ACBA},\qquad v'=\mathtt{BAAC},
\]
which are again members of $\mathcal F$. Unlike the previous example, the
resulting words are different from the original two words.
\end{example}

\begin{example}[A four-color cycle]
The family
\[
 \mathcal F=\{\mathtt{ABCD},\mathtt{BCDA}\}
\]
satisfies symmetric cycle exchange. Between its two words, all four positions
form the simple color cycle
\[
 A\longrightarrow B\longrightarrow C\longrightarrow D\longrightarrow A.
\]
The only possible switch interchanges the two words, and its cycle contains
every prescribed disagreement.
\end{example}

\section[twoColors]{The two-color case}

Recall that a nonempty equicardinal family of subsets of $E$ is the family of
bases of a matroid if it satisfies basis exchange: for bases $B,C$ and
$e\in B\setminus C$, there is an $f\in C\setminus B$ such that
$B-e+f$ is a basis. We write $\mathcal B(M)$ for the bases of a matroid $M$.
The dual matroid $M^*$ has bases
\[
 \mathcal B(M^*)=\{E\setminus B:B\in\mathcal B(M)\}.
\]

\begin{proposition}
Let $\mathcal F$ be a fixed-quota coloring family with colors $A$ and $B$.
Then $\mathcal F$ is an SCE family if and only if there is a matroid $M$ on
$E$ such that
\[
 \mathcal F
 =\{u:E\to\{A,B\}:B(u,A)\in\mathcal B(M)\}.
\]
Under this correspondence, the color-$A$ classes are the bases of $M$, and
the color-$B$ classes are the bases of $M^*$.
\end{proposition}

\begin{proof}
Suppose first that $\mathcal F$ is an SCE family. Take $u,v\in\mathcal F$ and
$e\in B(u,A)\setminus B(v,A)$. With two colors, a simple color cycle through
the arc at $e$ must be a two-cycle. Its other arc is labeled by some
$f\in B(v,A)\setminus B(u,A)$. Switching the two positions gives a word in
$\mathcal F$ whose color-$A$ class is
\[
 B(u,A)-e+f.
\]
Thus the color-$A$ classes satisfy basis exchange. They all have cardinality
$r_A$, so they are the bases of a matroid $M$. Since
$B(u,B)=E\setminus B(u,A)$, the color-$B$ classes are the bases of $M^*$.

Conversely, let $M$ be a matroid on $E$ and form $\mathcal F$ by assigning
color $A$ to the elements of a basis and color $B$ to the remaining elements.
Take
$u,v\in\mathcal F$ and a position $e$ where they disagree. After interchanging
$u$ and $v$ if needed, we may assume that
$e\in B(u,A)\setminus B(v,A)$. The symmetric basis-exchange theorem gives an
$f\in B(v,A)\setminus B(u,A)$ such that both
\[
 B(u,A)-e+f
 \qquad\text{and}\qquad
 B(v,A)-f+e
\]
are bases of $M$. The positions $e,f$ form a simple color two-cycle, and
switching them in both words gives two words in $\mathcal F$. Hence
$\mathcal F$ is an SCE family.
\end{proof}
